\subsection{Три способа организовать программу}

\subsubsection{1. Последовательность с sleep}
Действие, пауза, действие. Просто, но \hyperlink{term:blocking}{блокирует} программу: в полёте недопустимо. Годится только для проверки на земле.

\subsubsection{2. Таймеры}
Каждая анимация --- свой таймер, каждый тик --- один шаг. Таймеры независимы: мигание навигационных огней и перелив матрицы работают одновременно, а программа остаётся свободной для событий.

\subsubsection{3. Конечный автомат}
Когда поведение зависит от того, «что уже произошло», одной переменной-флага мало. \hyperlink{term:automaton}{Конечный автомат} --- это:
\begin{itemize}
    \item множество \hyperlink{term:state}{состояний} (\verb|"RED"|, \verb|"YELLOW"|, \verb|"GREEN"|);
    \item события (тик таймера, \verb|Ev.POINT_REACHED|, тумблер пульта);
    \item правила переходов: «в состоянии A пришло событие X $\to$ перейти в B и сделать Y».
\end{itemize}

Светофор --- самый маленький автомат:

\begin{codefile}{lua/examples/05_leds/leds_traffic_light.lua}
\end{codefile}

\paragraph{Что происходит.} Таблица \verb|nextState| --- это и есть правила переходов, \verb|colors| --- действие в каждом состоянии. Таймер приносит событие «прошло 2 секунды», обработчик делает переход и действие. Добавить четвёртое состояние --- значит добавить по строке в две таблицы, не трогая логику.

\subsection{Примеры комплексных сценариев}

\subsubsection*{Перелив цвета и аварийная индикация}

Этот скрипт --- официальный пример Геоскан, слегка причёсанный:

\begin{codefile}{lua/examples/05_leds/leds_random_color_example.lua}
\end{codefile}

Таймер каждые 0.05~с чуть сдвигает одну из компонент цвета --- получается плавный перелив. Событие \verb|LOW_VOLTAGE2| через функцию \verb|emergency| останавливает таймер и включает красный. Два механизма --- таймер и событие --- работают вместе, не мешая друг другу.

\subsubsection*{Индикация этапов полёта}

А это --- мост к следующей главе. Программа взлетает, летит к точке и садится, а цвет линейки показывает этап:

\begin{codefile}{lua/examples/05_leds/leds_mission_stages.lua}
\end{codefile}

Белый --- подготовка, зелёный --- полёт к точке, синий --- посадка, выключено --- дрон на земле. Команды \verb|ap.push| и \verb|ap.goToLocalPoint| разобраны в главе~\ref{sec:missions}; здесь важно другое: каждая смена цвета стоит рядом с переходом на новый этап, и по дрону видно, где программа. Переменной \verb|state| в этой программе нет: этапы идут строго друг за другом, и каждое событие (\verb|TAKEOFF_COMPLETE|, \verb|POINT_REACHED|, \verb|COPTER_LANDED|) приходит ровно один раз, поэтому само событие однозначно говорит, какой этап начался. Когда одно и то же событие может прийти на разных этапах, понадобится явная переменная состояния --- так устроены миссии в разделе~\ref{sec:missions-architecture}.
